\section{Proje Planı}
\label{BesinciBolum}

Projenin kilometre taşları ve 3 Kasım 2026 itibarıyla durumları
Tablo~\ref{tab:plan} içinde verilmiştir. M1 kapsamındaki nokta çıkarımı, BiGRU ve
1B-CNN + Transformer modellerinin 60 epoch eğitimi, hata analizi ve LiteRT
dönüşümü tamamlanmış; çoklu tohum (seed) koşuları ve mobil uygulamanın kalan
ekranları M3 ve M4 altında planlanmıştır.

\begin{table}[H]
\centering
\caption{Kilometre taşları ve durumları}
\label{tab:plan}
\small
\begin{tabularx}{\textwidth}{@{}l l l X@{}}
\toprule
\textbf{Kilometre Taşı} & \textbf{Tarih} & \textbf{Durum} & \textbf{Çıktı / Açıklama} \\
\midrule
T1 & 5 Ekim 2026 & Tamamlandı & Açılış toplantısı ve iş bölümü. \\
M0 & 29 Ekim 2026 & Tamamlandı & Depo ve erişim izinleri, AUTSL verisinin doğrulanması. \\
M1 & 7 Kasım 2026 & Tamamlandı & Nokta çıkarımı, BiGRU ve 1B-CNN+Transformer eğitimi, fp32/int8 LiteRT modelleri. \\
M2 & 20 Kasım 2026 & Devam ediyor & Ara raporun ve model karşılaştırma/hata analizi grafiklerinin teslimi. \\
M3 & 4 Aralık 2026 & Planlandı & Çoklu tohum (3 seed) standart sapma analizi ve nihai model dondurma. \\
M4 & 11 Aralık 2026 & Planlandı & Üç ekranlı mobil uygulamanın ve uçak modu testlerinin tamamlanması. \\
M5 & 18 Aralık 2026 & Planlandı & Gönüllü testleri ve performans ölçümlerinin tamamlanması. \\
M6 & 1 Ocak 2027 & Planlandı & Tez teslimi ve Turnitin kontrolü. \\
\bottomrule
\end{tabularx}
\end{table}

Ara sınav haftası (9--15 Kasım) kod yazmak yerine her iki modelin üç farklı
tohumla (\texttt{seed 0, 1, 2}) arka planda koşulmasına ayrılmıştır; böylece
sınav dönüşünde ortalama ve standart sapma tabloları hazır olacaktır.
